% =============================================================
%  附录 A —— 物理量符号及含义
%  按正文首次出现的先后排列；**不记页码**，改为给出每个物理量的量纲
%  （用户拍板，第 5 章交付后）。
% =============================================================
\bichapter{物理量符号及含义}{List of Symbols}

本附录按正文中首次出现的先后顺序（按章分组），列出全书出现的物理量与关键统计量符号、含义及其量纲。量纲一栏以基本量纲 $\mathsf{M}$（质量）、$\mathsf{L}$（长度）、$\mathsf{T}$（时间）、$\Theta$（温度）、$\mathsf{I}$（电流）以及粒子数 $\mathsf{N}$ 表示；能量一律记作 $\mathsf{M}\mathsf{L}^2\mathsf{T}^{-2}$，动量记作 $\mathsf{M}\mathsf{L}\mathsf{T}^{-1}$，纯计数或无单位的概率、比例标为「无量纲」。同一符号在不同章节表示不同物理量时（如 $\Omega$、$\rho$、$\mu$），在相应条目中分别给出。

\begin{xltabular}{\textwidth}{@{}l X l@{}}
  \toprule
  \textbf{符号} & \textbf{含义} & \textbf{量纲} \\
  \midrule
  \endfirsthead
  \toprule
  \textbf{符号} & \textbf{含义} & \textbf{量纲} \\
  \midrule
  \endhead
  \midrule
  \multicolumn{3}{r@{}}{\small 续下页} \\
  \endfoot
  \bottomrule
  \endlastfoot

  \multicolumn{3}{@{}l}{\textbf{第 1 章\hspace{0.4em}热力学}} \\
  \midrule
  $T$ & 绝对温度 & $\Theta$ \\
  $\Theta$ & 经验温度（温标） & $\Theta$ \\
  $\inexd Q$ & 系统吸收的微元热量 & $\mathsf{M}\mathsf{L}^2\mathsf{T}^{-2}$ \\
  $\inexd W$ & 系统对外做的微元功 & $\mathsf{M}\mathsf{L}^2\mathsf{T}^{-2}$ \\
  $U$ & 内能 & $\mathsf{M}\mathsf{L}^2\mathsf{T}^{-2}$ \\
  $P$ & 压强 & $\mathsf{M}\mathsf{L}^{-1}\mathsf{T}^{-2}$ \\
  $V$ & 体积 & $\mathsf{L}^3$ \\
  $N$ & 粒子数（全书亦用作粒子数的通记） & 无量纲（计数） \\
  $n=N/V$ & 数密度 & $\mathsf{L}^{-3}$ \\
  $k_B$ & Boltzmann 常量 & $\mathsf{M}\mathsf{L}^2\mathsf{T}^{-2}\Theta^{-1}$ \\
  $S$ & 熵 & $\mathsf{M}\mathsf{L}^2\mathsf{T}^{-2}\Theta^{-1}$ \\
  $C$ & 热容 & $\mathsf{M}\mathsf{L}^2\mathsf{T}^{-2}\Theta^{-1}$ \\
  $c$ & 比热容（单位质量） & $\mathsf{L}^2\mathsf{T}^{-2}\Theta^{-1}$ \\
  $L$ & 潜热（每分子） & $\mathsf{M}\mathsf{L}^2\mathsf{T}^{-2}$ \\
  $\mu$ & 化学势（每增加一个粒子所增加的能量） & $\mathsf{M}\mathsf{L}^2\mathsf{T}^{-2}$ \\
  $H=U+PV$ & 焓 & $\mathsf{M}\mathsf{L}^2\mathsf{T}^{-2}$ \\
  $F=U-TS$ & Helmholtz 自由能 & $\mathsf{M}\mathsf{L}^2\mathsf{T}^{-2}$ \\
  $\mathcal{G}=U-TS+PV$ & Gibbs 自由能 & $\mathsf{M}\mathsf{L}^2\mathsf{T}^{-2}$ \\
  $\kappa_T=-\frac1V\left.\frac{\partial V}{\partial P}\right|_T$ & 等温压缩率 & $\mathsf{M}^{-1}\mathsf{L}\mathsf{T}^{2}$ \\
  $\alpha=\frac1V\left.\frac{\partial V}{\partial T}\right|_P$ & 等压体膨胀系数 & $\Theta^{-1}$ \\
  $\gamma=C_P/C_V$ & 绝热指数（比热比） & 无量纲 \\
  $R$ & 气体常量（每 mol） & $\mathsf{M}\mathsf{L}^2\mathsf{T}^{-2}\Theta^{-1}\mathsf{N}^{-1}$ \\
  $\eta_{\text{Carnot}}$ & Carnot 效率 & 无量纲 \\
  $\mu_B$ & Bohr 磁子（微观磁矩） & $\mathsf{I}\mathsf{L}^2$ \\

  \multicolumn{3}{@{}l}{\textbf{第 2 章\hspace{0.4em}概率论}} \\
  \midrule
  $\mathscr{S}$ & 结果（样本）集合 & 无量纲（集合） \\
  $E$ & 事件（结果集合的任意子集） & 无量纲（集合） \\
  $p(x)$ & 概率密度函数（PDF，连续变量） & $[x]^{-1}$ \\
  $P(E)$ & 事件 $E$ 的概率 & 无量纲 \\
  $\langle X\rangle$ & 随机变量 $X$ 的期望值 & $[X]$ \\
  $\sigma^2$ & 方差 & $[X]^2$ \\
  $\sigma$ & 标准差 & $[X]$ \\
  $m_n=\langle X^n\rangle$ & 第 $n$ 阶矩 & $[X]^n$ \\
  $\langle\langle X^n\rangle\rangle$ & 第 $n$ 阶累积量 & $[X]^n$ \\
  $\tilde p(k)$ & 特征函数（PDF 的 Fourier 变换） & 无量纲 \\
  $\log\tilde p(k)$ & 累积量生成函数（$\log$ 特征函数） & 无量纲 \\
  $\mathscr{N}$ & 归一化常量 & $[x]^{-1}$ \\

  \multicolumn{3}{@{}l}{\textbf{第 3 章\hspace{0.4em}气体动理论}} \\
  \midrule
  $f_r(\vec p,\vec q,t)$ & $r$ 粒子分布函数（相空间中的粒子数密度） & $\mathsf{M}^{-3}\mathsf{L}^{-6}\mathsf{T}^{3}$ \\
  $n(\vec q,t)$ & 局域粒子数密度 & $\mathsf{L}^{-3}$ \\
  $\rho(\bm p,\bm q,t)$ & 相空间中的概率密度（按 $\mathrm d\Gamma$ 测度归一） & $\mathsf{M}^{-3}\mathsf{L}^{-6}\mathsf{T}^{3}$ \\
  $\mathscr H$ & Hamiltonian（系统的总能量） & $\mathsf{M}\mathsf{L}^2\mathsf{T}^{-2}$ \\
  $\mathscr V$ & 粒子间的两体相互作用势 & $\mathsf{M}\mathsf{L}^2\mathsf{T}^{-2}$ \\
  $\vec b$ & 碰撞参数 & $\mathsf{L}$ \\
  $\sigma$ & 散射截面 & $\mathsf{L}^2$ \\
  $\mathrm{d}\sigma/\mathrm{d}\Omega$ & 微分散射截面 & $\mathsf{L}^2$ \\
  $\Omega$ & 立体角（散射问题中的末态方向） & 无量纲 \\
  $\ell$ & 平均自由程 & $\mathsf{L}$ \\
  $\tau_\times$ & 两次碰撞之间的平均自由时间 & $\mathsf{T}$ \\
  $\tau_c$ & 一次碰撞的持续时间 & $\mathsf{T}$ \\
  $\bar n,\ \bar T$ & 局域平衡的数密度与温度 & $\mathsf{L}^{-3}$、$\Theta$ \\
  $\nu$ & 运动粘度（动量扩散系数） & $\mathsf{L}^2\mathsf{T}^{-1}$ \\
  $K$ & 热导率 & $\mathsf{M}\mathsf{L}\mathsf{T}^{-3}\Theta^{-1}$ \\
  $\vec J_n=n\vec u$ & 粒子流密度 & $\mathsf{L}^{-2}\mathsf{T}^{-1}$ \\
  $J_\chi$ & 守恒量 $\chi$ 的碰撞产生率 & 视 $\chi$ 而定 \\
  $v$ & 纵模的声速 & $\mathsf{L}\mathsf{T}^{-1}$ \\
  $\tilde u,\tilde\theta,\tilde\nu$ & 速度、温度与密度的微扰 & $\mathsf{L}\mathsf{T}^{-1}$、$\Theta$、$\mathsf{L}^{-3}$ \\

  \multicolumn{3}{@{}l}{\textbf{第 4 章\hspace{0.4em}经典统计力学}} \\
  \midrule
  $\mu$ & 相空间中的一个点（第 4 章的记号） & --- \\
  $\mathbf x$ & 宏观态中除能量外的其它参量（如总磁化） & 视参量而定 \\
  $p(\mu)$ & 相空间上的概率密度 & 无量纲（归一） \\
  $\Omega(E,\mathbf x)$ & 能量 $E$ 处的可及相空间体积（微观态数目） & 无量纲 \\
  $Z,\ \mathcal{Z}$ & 正则配分函数 & 无量纲 \\
  $\mathcal{Q}$ & 巨配分函数 & 无量纲 \\
  $\mathcal{G}$ & 巨势（巨正则自由能） & $\mathsf{M}\mathsf{L}^2\mathsf{T}^{-2}$ \\
  $\beta=1/k_BT$ & 逆温度 & $\mathsf{M}^{-1}\mathsf{L}^{-2}\mathsf{T}^{2}$ \\
  $\mathcal{E}$ & 能量（随机变量） & $\mathsf{M}\mathsf{L}^2\mathsf{T}^{-2}$ \\
  $\lambda=h/\sqrt{2\pi mk_BT}$ & 热波长 & $\mathsf{L}$ \\
  $p(\mathcal{E})$ & 能量的概率分布 & $(\mathsf{M}\mathsf{L}^2\mathsf{T}^{-2})^{-1}$ \\
  $M=\mu_0\sum_i\sigma_i$ & 系统的总磁化 & $\mathsf{I}\mathsf{L}^2$ \\
  $\mu_0$ & 单个自旋的微观磁矩 & $\mathsf{I}\mathsf{L}^2$ \\
  $B$ & 外加磁场 & $\mathsf{M}\mathsf{T}^{-2}\mathsf{I}^{-1}$ \\
  $\chi(T)=\left.\partial M/\partial B\right|_{B=0}$ & 磁化率 & $\mathsf{I}^2\mathsf{L}^2\mathsf{M}^{-1}\mathsf{T}^{2}$ \\

  \multicolumn{3}{@{}l}{\textbf{第 5 章\hspace{0.4em}相互作用粒子}} \\
  \midrule
  $\mathcal U$ & 粒子的相互作用总势能 & $\mathsf{M}\mathsf{L}^2\mathsf{T}^{-2}$ \\
  $u(r)$ & 两体势（$u_0$ 为其极小值深度） & $\mathsf{M}\mathsf{L}^2\mathsf{T}^{-2}$ \\
  $f(\vec q)=\mathrm e^{-\beta\mathscr V}-1$ & Mayer $f$ 函数 & 无量纲 \\
  $\mathscr{S}_N$ & $N$ 粒子集团积分 & $\mathsf{L}^{3(N-1)}$ \\
  $\bar b_\ell$ & 不可约集团积分 & $\mathsf{L}^{3(\ell-1)}$ \\
  $B_i(T)$ & 第 $i$ 维里系数 & $\mathsf{L}^{3(i-1)}$ \\
  $x=\mathrm e^{\beta\mu}/\lambda^3$ & 逸度（活度） & $\mathsf{L}^{-3}$ \\
  $Q(\mu,T,V)$ & 巨配分函数（第 5 章记号） & 无量纲 \\
  $\rho_{\text{liquid}},\rho_{\text{gas}}$ & 液相与气相的密度 & $\mathsf{M}\mathsf{L}^{-3}$ \\
  $T_c,\ P_c,\ \rho_c$ & 临界温度、临界压强与临界密度 & $\Theta$、$\mathsf{M}\mathsf{L}^{-1}\mathsf{T}^{-2}$、$\mathsf{M}\mathsf{L}^{-3}$ \\
  $\alpha,\beta,\gamma,\delta,\nu$ & 临界指数 & 无量纲 \\

  \multicolumn{3}{@{}l}{\textbf{第 6 章\hspace{0.4em}量子统计力学}} \\
  \midrule
  $q_{i,\alpha}$ & 离子偏离平衡位置的坐标 & $\mathsf{L}$ \\
  $K_s$ & 简正模的劲度（弹簧常量） & $\mathsf{M}\mathsf{T}^{-2}$ \\
  $\mathcal V(\vec q)$ & 势能（离子构型的函数） & $\mathsf{M}\mathsf{L}^2\mathsf{T}^{-2}$ \\
  $\vec k$ & 波矢 & $\mathsf{L}^{-1}$ \\
  $a$ & 晶格常量 & $\mathsf{L}$ \\
  $\omega$ & 振动模式的角频率 & $\mathsf{T}^{-1}$ \\
  $T_D$ & 德拜温度 & $\Theta$ \\
  $\sigma$ & Stefan--Boltzmann 常量 & $\mathsf{M}\mathsf{T}^{-3}\Theta^{-4}$ \\
  $\hbar$ & 约化 Planck 常量 & $\mathsf{M}\mathsf{L}^2\mathsf{T}^{-1}$ \\
  $\rho$ & 密度矩阵 & 无量纲 \\
  $\Psi$ & 态向量（波函数） & 无量纲 \\
  $\mathcal E_n$ & 单粒子能量的本征值 & $\mathsf{M}\mathsf{L}^2\mathsf{T}^{-2}$ \\
  $\mathcal N$ & 微观态的数目 & 无量纲（计数） \\
  $p_\alpha$ & 系统处于微观态 $\alpha$ 的概率 & 无量纲 \\

  \multicolumn{3}{@{}l}{\textbf{第 7 章\hspace{0.4em}理想量子气体}} \\
  \midrule
  $\eta$ & 统计符号（玻色子 $+1$、费米子 $-1$） & 无量纲 \\
  $s$ & 自旋量子数（以 $\hbar$ 为单位） & 无量纲 \\
  $g=2s+1$ & 自旋简并度 & 无量纲 \\
  $\mathcal E(\vec k)$ & 单粒子能量（波矢的函数） & $\mathsf{M}\mathsf{L}^2\mathsf{T}^{-2}$ \\
  $z=\mathrm e^{\beta\mu}$ & 逸度（第 7 章的记号） & 无量纲 \\
  $f^\pm_\alpha(z)$ & 费米（$+$）与玻色（$-$）积分 & 无量纲 \\
  $n_\eta$ & 粒子数密度（$\eta$ 标记统计） & $\mathsf{L}^{-3}$ \\
  $k_F$ & 费米波数 & $\mathsf{L}^{-1}$ \\
  $\mathcal E_F$ & 费米能 & $\mathsf{M}\mathsf{L}^2\mathsf{T}^{-2}$ \\
  $T_F=\mathcal E_F/k_B$ & 费米温度 & $\Theta$ \\
  $P_F=(2/5)n\mathcal E_F$ & 费米压强 & $\mathsf{M}\mathsf{L}^{-1}\mathsf{T}^{-2}$ \\
  $\zeta_{3/2},\ \zeta_{5/2}$ & Riemann $\zeta$ 函数在 $3/2$、$5/2$ 处的值 & 无量纲 \\
  $T_c$ & 玻色--爱因斯坦凝聚温度 & $\Theta$ \\
  $n_0$ & 凝聚于 $\vec k=0$ 态的粒子数密度 & $\mathsf{L}^{-3}$ \\
  $n^*$ & 激发态的极限粒子数密度 & $\mathsf{L}^{-3}$ \\
  $\rho_n,\ \rho_s$ & 正常成分与超流成分的质量密度 & $\mathsf{M}\mathsf{L}^{-3}$ \\
  $s_n$ & 正常成分的熵密度 & $\mathsf{M}\mathsf{L}^{-1}\mathsf{T}^{-2}\Theta^{-1}$ \\
  $\vec v_n,\ \vec u_s$ & 正常成分与超流成分的速度 & $\mathsf{L}\mathsf{T}^{-1}$ \\
  $\Delta$ & 旋子谱的能隙 & $\mathsf{M}\mathsf{L}^2\mathsf{T}^{-2}$ \\
  $k_0$ & 旋子谱极小处的波数 & $\mathsf{L}^{-1}$ \\
\end{xltabular}
